%can_dphi_fits_cent0_Lorentz_EffCor

\subsubsection{Fits to the \Dphi\ distributions}

In order to obtain the \iftoggle{IncludeYields}{yield and}{} widths of the distributions
  under the away-side peak, the \Dphi\ distributions are fit 
  with the following functional form:
\begin{eqnarray}
dN/d\Dphi=N_{0}(1+v^{\mathrm{flow}}_{2,2}\cos(2\Dphi))+ Y^{\mathrm{corr}}(\Dphi),
\label{eq:fit_func}
\end{eqnarray}
with:
\begin{eqnarray}
%Y^{\mathrm{corr}} (\Dphi) = N^{\mathrm{corr}}\frac{\Gamma/\pi}{(\Dphi-\pi)^2 + \Gamma^2}- N^{\mathrm{pedestal}},
Y^{\mathrm{corr}} (\Dphi) = \frac{N^{\mathrm{corr}}}{(\Dphi-\pi)^2 + \Gamma^2}- N^{\mathrm{pedestal}},
\label{eq:signal_func}
\end{eqnarray}
where $N_0$, $v^{\mathrm{flow}}_{2,2}$, $N^{\mathrm{corr}}$ and $\Gamma$ are fit-parameters.
The $N^{\mathrm{pedestal}}$ term is constrained to ensure that 
  $Y^{\mathrm{corr}} (\Dphi)$ is zero at $\Dphi=0$.
The explanation of the terms in the fit is as follows.
\begin{itemize}
  \item $N_0$ : 
    This term is independent of \Dphi\ and is used to account for
      the UE pairs that have no modulation in \Dphi.

  \item $v^{\mathrm{flow}}_{2,2}\cos(2\Dphi)$ : 
    This term accounts for the (elliptic or $2^{\mathrm{nd}}$-order) flow modulation in the UE pairs.

  \item $Y^{\mathrm{corr}}(\Dphi)$ : 
    This term parameterizes the correlated yield.
    The yield is parameterized as a Lorentzian function, 
      centered at $\Dphi=\pi$, with a pedestal that is adjusted
      to ensure $Y^{\mathrm{corr}} (\Dphi)=0$ at $\Dphi=0$.
    The Lorentzian form is chosen as it describes the data well.
    The $Y^{\mathrm{corr}}(\Dphi)$ is folded over the range 
      ($-\pi/2,3\pi/2$).
    The integral of $Y^{\mathrm{corr}}(\Dphi)$, over the full \Dphi\ range,
      is the correlated yield.
    The standard-deviation of $Y^{\mathrm{corr}}(\Dphi)$, 
      (and not the parameter $\Gamma$) is taken as
      the width of the correlation.
    While calculating the standard-deviation of $Y^{\mathrm{corr}}(\Dphi)$
      its folding into the range ($-\pi/2,3\pi/2$) is taken into consideration.
    This is achieved by remapping the $Y^{\mathrm{corr}}(\Dphi)$
      over the range ($0,2\pi$), so that the away-side peak is centered, 
      and then calculating its RMS-width.
    \iftoggle{IncludeYields}
    {
      The correlated yields and the RMS-widths are calculated directly 
        using the fit function Eq.~\ref{eq:signal_func}.
      Statistical uncertainties in the yields and RMS-widths are
        obtained using toys, where $dN/d\Dphi$ distribution
        are resampled and refitted using Eq.\eqref{eq:fit_func},
        and the standard deviations in the yields and RMS-widths
        in the resamplings are taken as the corresponding uncertainties.
    }
    {
      The RMS-widths of the correlated pair-distribution are calculated directly 
        using the fit function Eq.~\ref{eq:signal_func}.
      Statistical uncertainties of the RMS-widths are
        obtained using toys, where $dN/d\Dphi$ distribution
        are resampled and refitted using Eq.\eqref{eq:fit_func},
        and the standard deviations of the RMS-widths
        in the resamplings are taken as the corresponding uncertainty.
    }
    Plots for the uncertainty estimation using toys are included in Appendix~\ref{sec:stat_errs}.
\end{itemize}


Figure~\ref{fig:fit_PbPb_tight} shows fits to the \Dphi\ distributions
  using Eq.~\eqref{eq:fit_func}, for the \PbPb\ data.
The fits are shown for same-charge, opposite-charge pairs as well 
  as the combined case, and for several centrality intervals.
And similar plots for the \pp\ data are shown in Figure~\ref{fig:fit_pp}.
For completeness, Figure~\ref{fig:v22_centdep0} and Figure~\ref{fig:tau_centdep0}
 in Appendix~\ref{sec:fit_parameters} show the fit-parameters $v_{2,2}^{flow}$ and $\Gamma$
 from these fits.
%In all cases, the fits describe the data well.



\begin{figure}[h]
\centering
\includegraphics[width=1.0\textwidth]%
{figures/AllFigs/can_dphi_fits_cent12_LorentzZYAM_Tight_pTBar_GT5_EffCor}
\includegraphics[width=1.0\textwidth]%
{figures/AllFigs/can_dphi_fits_cent13_LorentzZYAM_Tight_pTBar_GT5_EffCor}
\includegraphics[width=1.0\textwidth]%
{figures/AllFigs/can_dphi_fits_cent4_LorentzZYAM_Tight_pTBar_GT5_EffCor}
\includegraphics[width=1.0\textwidth]%
{figures/AllFigs/can_dphi_fits_cent6_LorentzZYAM_Tight_pTBar_GT5_EffCor}
\includegraphics[width=1.0\textwidth]%
{figures/AllFigs/can_dphi_fits_cent8_LorentzZYAM_Tight_pTBar_GT5_EffCor}
\caption{
Fits to the \Dphi\ distributions using Eq.~\eqref{eq:fit_func}.
Plots are for the \PbPb\ data.
The left, middle and right panels correspond to same-charge, 
  opposite-charge and combined-charge pairs, respectively. 
From top to bottom, each row corresponds to a different centrality interval.
Note: the $y$-range of the different panels are very different,
  and they are zero suppressed.
\label{fig:fit_PbPb_tight}
}
\end{figure}

\begin{figure}[h]
\centering
\includegraphics[width=1.0\textwidth]%
{figures/AllFigs/can_dphi_fits_cent11_LorentzZYAM_pTBar_GT5_pp2017_EffCor}
\includegraphics[width=1.0\textwidth]%
{figures/AllFigs/can_dphi_fits_cent11_LorentzZYAM_Tight_pTBar_GT5_pp2017_EffCor}
\caption{
Similar to Figure~\ref{fig:fit_PbPb_medium}, but for the \pp \ data.
The top and bottom panels correspond to Medium and Tight muons, respectively.
\label{fig:fit_pp}
}
\end{figure}






\clearpage
\FloatBarrier
\iftoggle{IncludeYields}
{
  \subsubsection{Correlated yields and widths}
}
{
  \subsubsection{Widths for the correlated muon-pair distributions}
}

Figure~\ref{fig:widths_PbPb} shows the widths for the \PbPb\ data 
  (standard deviation of $Y^{\mathrm{corr}}(\Dphi)$), 
  plotted as a function of centrality.
The results are shown for both the Tight and Medium selections.
The \PbPb\ data are also compared to the \pp\ measurements 
  (continuous line).
The widths for the \PbPb\ show no systematic dependence on centrality
  over the mid-central and peripheral range, where the widths are 
  consistent with the widths measured in the \pp\ data.
However, over the 0--20\% most-central events the same-charge pairs
  do show a slight narrowing in the widths compared to the more 
  peripheral intervals.

\begin{figure}[h]
\centering
\includegraphics[width=1.0\textwidth]%
{figures/AllFigs/can_widths_pTBar_GT5_EffCor}
\includegraphics[width=1.0\textwidth]%
{figures/AllFigs/can_widths_Tight_pTBar_GT5_EffCor}
\caption{
Widths of the away-side correlation in the \PbPb\ data 
  as a function of centrality.
The measurements are shown for the same-charge (left),
  opposite-charge (middle), and combined-charge (right) pairs.
The top and bottom rows correspond to Medium and Tight muons, respectively.
The \pp\ measurements are also shown (blue line). 
\label{fig:widths_PbPb}
}
\end{figure}




\iftoggle{IncludeYields}
{
  Figure~\ref{fig:yields_PbPb} shows the yields of the correlated pairs
    in different centrality intervals.
  The yields increase systematically from peripheral to central intervals.
  The number of opposite-sign pairs are approximately double the number 
    of same-sign pairs.

  Figure~\ref{fig:scaled_yields_PbPb} shows the per-event yields scaled by the nuclear 
    overlap function \TAA\ that is obtained using a Glauber model analysis~\cite{Miller:2007ri}:
  \begin{eqnarray}
  \text{Scaled per-event yield}=\frac{Yield}{N_{\mathrm{events}}^{\mathrm{PbPb}}\TAA},
  \label{eq:TaaYield}
  \end{eqnarray}
  where \textit{Yield} is the measured yield in a given 
    centrality interval (Figure~\ref{fig:yields_PbPb}),
    and $N_{\mathrm{events}}^{\mathrm{PbPb}}$ is the number
    of \PbPb\ events in the given centrality interval
    corresponding to the luminosity of 1.93\inb 
    (which can be obtained using $\sigma^{\PbPb}=7.66\mu b$~\cite{Loizides:2017ack}).
  Table~\ref{tbl:TAA} lists the \TAA\ values provided by the Heavy-ion group
    (\href{https://twiki.cern.ch/twiki/pub/AtlasProtected/HeavyIonAnalysis2015/centrality_values_Gv32_update_26Jul19.txt}{link})
    for the centrality intervals used in this analysis.
  The \TAA\ scaled yields (which are equivaled to a cross-section) decrease 
    considerably from peripheral to central colisions: more than a factor 
    of 5 decrease from the 60--80\% centrality interval to the 
    0--10\% centrality interval.
  This implies that a significant increase in the suppression of correlated 
    muon pairs is observed from peripheral to central collisions.
  Figure~\ref{fig:scaled_yields_PbPb} also shows the measured cross-sections 
    for the \pp\ data.
  The \TAA\ scaled yields in \PbPb\ are smaller than the \pp\ cross-sections
    even in the 60--80\% central events, indicating some suppression
    in this peripheral interval.


  \begin{table}[h]
  \centering
  \begin{tabular}{|c|c|c|}
  \hline
  Centrality [\%] & $T_{AA}$ [mb$^{-1}$]  & $\Delta T_{AA}$ [mb$^{-1}$] \\
  \hline
   0--10 & 23.2  & 0.124\\
  10--20 & 14.4  & 0.132\\
  20--30 &  8.77 & 0.132\\
  30--40 &  5.09 & 0.120\\
  40--50 &  2.75 & 0.095\\
  50--60 &  1.35 & 0.065\\
  60--80 &  0.42 & 0.003\\
  \hline
  \end{tabular}
  \caption{
  The \TAA\ values for the different centrality intervals,
    and the corresponding uncertainties.
  \label{tbl:TAA}
  }
  \end{table}






  \begin{figure}[h]
  \centering
  \includegraphics[width=1.0\textwidth]%
  {figures/AllFigs/can_yields_option1_pTBar_GT5_EffCor}
  \includegraphics[width=1.0\textwidth]%
  {figures/AllFigs/can_yields_option1_Tight_pTBar_GT5_EffCor}
  \caption{
  Yields of the correlated muon-pairs on the away-side 
    in the \PbPb\ data, as a function of centrality.
  The measurements are shown for the same-charge (left),
    opposite-charge (middle), and combined-charge (right) pairs.
  The top and bottom rows correspond to Medium and Tight muons, respectively.
  \label{fig:yields_PbPb}
  }
  \end{figure}


  \begin{figure}[h]
  \centering
  \includegraphics[width=1.0\textwidth]%
  {figures/AllFigs/can_yields_option2_pTBar_GT5_EffCor}
  \includegraphics[width=1.0\textwidth]%
  {figures/AllFigs/can_yields_option2_Tight_pTBar_GT5_EffCor}
  \caption{
  \TAA\ scaled per-event Yields of the correlated muon-pairs on the away-side 
    in the \PbPb\ data, as a function of centrality.
  The measurements are shown for the same-charge (left), 
    opposite-charge (middle), and combined-charge (right) pairs.
  The top and bottom rows correspond to Medium and Tight muons, respectively.
  The cross-section in the \pp\ measurements are also shown (blue line). 
  \label{fig:scaled_yields_PbPb}
  }
  \end{figure}
}
{
  
}
